\chapter{Configurations and Run Commands}
\label{app:hyperparams}


% Figure and table captions in this appendix are set at footnotesize, single
% spaced, as they were when these figures lived in their own appendix; moving
% them must not change how they print.
\captionsetup{font={footnotesize,singlespacing}}
Chapter~\ref{ch:data} states the experimental design in prose; this appendix gives
the fields, so that a reader can reconstruct any arm rather than take the design
on trust. Every value below carries a source comment naming where it was
read. Most name the version-controlled configuration layer
(\texttt{configs/base.yaml}, \texttt{configs/models/*.yaml},
\texttt{configs/hpo\_arm.yaml}) or the released per-run \texttt{config.json}
preimages listed in Appendix~\ref{app:external}; the rest name the run store,
the evidence tables, or the code that fixes a semantic no configuration file
records. All can be checked mechanically against the repository. % src: the 88 source comments in this file, by first path: 28 results/tables/, 16 configs/, 14 results/runs/, 8 scripts/experiments/, 8 release/raw_generations/, 5 release/run_configs/, 9 elsewhere (src/, scripts/analysis/, results/hpo/, release/ documents)

\FloatBarrier
\section{Shared Defaults and the One Training Recipe}
\label{app:hyperparams:recipe}

All arms inherit \texttt{configs/base.yaml}: the pinned chronological splits
(train 2010--2019, validation 2020--2021, test 2022--2025, assigned by effective
trading day and day-inclusive on both boundaries), horizons $h \in \{5,10,20\}$
with metrics computed per horizon, the three disclosure subsets, a log
realised-volatility target, and default seed 2026.

One recipe, fixed before the production campaign, governs every fine-tuned
Block~C/D run: AdamW; learning rate $8\times10^{-5}$ (the pilot-era
$1\times10^{-5}$ linearly scaled to the unified effective batch); weight decay
0.01; effective batch \textbf{128}, reached physically or by gradient
accumulation; linear warmup over the first 6\% of steps then decay; bfloat16
mixed precision; at most 15 epochs with validation early stopping (patience 3,
minimum delta 0) and best-epoch restore; squared-error loss on the log target;
and a regression head of one hidden layer of width 128 with dropout 0.1. % src: configs/models/C*.yaml,D1,D2 training block; TECHNICAL_SUPPLEMENT.md Table S8(a)
Document text is pre-tokenised once and cached across horizons, epochs and seeds.

Frozen-embedding arms C5 and D3 use the head-only analogue: encoder weights
are never updated, each filing is encoded once and its embedding cached and
reused across seeds, horizons, splits and the C5/D3 pair, and only the head
trains. It runs at learning rate $1\times10^{-4}$, batch 512, epoch cap 40,
with the same warmup, early-stopping, precision and loss settings. C5x is the
parity arm and
trains nothing at all: a frozen encoder feeding a closed-form ridge head, so its
released configuration carries no training block and no training fingerprint. % src: configs/models/C5_*.yaml, D3_*.yaml
The larger cap is affordable precisely because the encoder is frozen.

That this is genuinely one recipe is machine-checkable rather than asserted: of
the 240 fingerprinted production runs, the 178 that carry a training block
collapse to 18 distinct training digests, clustered by model family, one of which
covers 53 of the 54 frozen-embedding runs; the odd one out,
\texttt{C5\_qwen3\_full\_long\_form\_seed2026}, carries a digest of its own. % src: results/tables/config_fingerprints.csv (54 run ids with prefix C5\_ or D3\_, all carrying a training digest, over two distinct values: 53 and 1)
Table~\ref{tab:app-settings} tabulates only the fields that vary.

\begin{table}[htbp]
  \centering
  \caption[Per-arm settings that vary across the fixed recipe]{Per-arm settings.
    Every field not shown is the shared recipe of
    Section~\ref{app:hyperparams:recipe}. ``Batch'' gives physical batch
    $\times$ accumulation steps. The four Block~B arms are ridge regressions and
    their configurations record no length, rate, batch or epoch field.
    Source: \texttt{configs/models/*.yaml}; \texttt{release/run\_configs/}.}
  \label{tab:app-settings}
  \footnotesize
  \setlength{\tabcolsep}{4pt}
  \begin{tabular}{@{}llrlr@{}}
    \toprule
    Arm & Checkpoint or class & Max len. & Batch & Epochs \\
    \midrule
    B1--B4 & ridge over lexical/dictionary features & --- & --- & --- \\
    \midrule
    C1 S1 & \texttt{bert-base-uncased} & 512 & $128\times1$ & 15 \\
    C1 S2 & \texttt{bert-base-uncased}, chunk mean & 512 & $8\times16$ & 15 \\
    C2 S1 & \texttt{ProsusAI/finbert} & 512 & $128\times1$ & 15 \\
    C2 S2 & \texttt{ProsusAI/finbert}, chunk mean & 512 & $8\times16$ & 15 \\
    C2 S3 & \texttt{ProsusAI/finbert}, chunk attention & 512 & $8\times16$ & 15 \\
    C2 S4 & \texttt{ProsusAI/finbert}, hierarchical & 512 & $8\times16$ & 15 \\
    C3 S1 & \texttt{roberta-base} & 512 & $128\times1$ & 15 \\
    C4 S5 & \texttt{allenai/longformer-base-4096} & 4{,}096 & $4\times32$ & 15 \\
    C5 & \texttt{Qwen/Qwen3-Embedding-8B} (frozen) & 4{,}096 & 512 & 40 \\
    C5 & \texttt{Alibaba-NLP/gte-Qwen2-7B-instruct} & 4{,}096 & 512 & 40 \\
    C5 & \texttt{intfloat/e5-mistral-7b-instruct} & 4{,}096 & 512 & 40 \\
    C5x & Qwen3-Embedding-8B on C6 excerpts & 4{,}096 & ridge head & --- \\
    C6 & Qwen3-32B, prompted, no training & 8{,}192 & --- & --- \\
    \midrule
    D1 & FinBERT $\oplus$ price, concatenation MLP & 512 & $128\times1$ & 15 \\
    D2 & FinBERT $\oplus$ price, gated & 512 & $128\times1$ & 15 \\
    D3 & each frozen C5 embedder $\oplus$ price & 4{,}096 & 512 & 40 \\
    D4 & Qwen3-32B, price lags in prompt & 8{,}192 & --- & --- \\
    \bottomrule
  \end{tabular}
\end{table}

Learning rates follow the block: $8\times10^{-5}$ for every fine-tuned arm of
the fixed recipe and $1\times10^{-4}$ for every head-only arm. The two
ASHA-tuned arms of Section~\ref{app:hyperparams:asha} depart from it, as that
section reports. % src: configs/models/*.yaml; TECHNICAL_SUPPLEMENT.md Table S8(a)

\FloatBarrier
\section{Per-Block Specifics}
\label{app:hyperparams:blocks}

\paragraph{Block~A.} A1 carries the trailing 22-day realised volatility forward
and has no fitted parameters. A2 HAR-RV is a \emph{pooled panel} OLS, one fit per
horizon within each disclosure channel rather than one per firm
(\texttt{fit\_per\_firm: false}), on the log of
realised volatility at lags 1, 5 and 22. Duan smearing on retransformation
makes the prediction target the mean rather than the median. % src: configs/models/A2_har_rv.yaml; results/runs/A2_har_rv_full_{long_form,event_driven,combined}_seed2026
A3 and A4 fit GARCH(1,1) and EGARCH(1,1) per firm on log returns up to filing
time, normal innovations, minimum history 100 observations; A5 fits ARIMA(1,0,1)
per firm on log trailing realised volatility, minimum history 60. % src: configs/models/A3_garch.yaml, A4_egarch.yaml, A5_arima.yaml
A6 carries two arms under A2's conventions: SHAR splits the daily lag into
signed semivariances, and HARQ shrinks the weight on that lag where it is
imprecisely measured. A7
adds the point-in-time log VIX (the last close strictly before the effective
trading day) to the same pooled log-space fit. % src: release/run_configs/A6_shar_*.json, A6_harq_*.json, A7_harx_vix_*.json

\paragraph{Why A2 is pooled.} A2 is the primary baseline for every significance
test in the study, so the decision to fit it across firms rather than firm by firm
deserves an argument and not only a configuration flag. The configuration is where
the reason was first written down: the flag sits beside the annotation that
per-firm estimation overfits short panels. % src: configs/models/A2_har_rv.yaml (fit_per_firm annotation)
Two concessions come before the argument. The field documents what the estimator
does rather than switching it. A per-firm-slope HAR was also never fitted: the
estimator implements one pooled fit, or one fit per horizon, so no per-firm
variant appears among the executed arms. % src: src/sp500vol/models/price/har_rv.py (HARRV.fit)
The four points below therefore bound that omission instead of closing it by
head-to-head comparison.

First, the panel is short per firm. The long-form training split holds 19{,}668
filings from 702 distinct filers, an average of 28 filings each, and 187 of those
filers contribute 16 filings or fewer. % src: release/split_definition.md (19,668 long-form train filings); release/accession_index.csv (distinct cik and per-filer counts, form 10-K/10-Q, split=train)
Because the fit is one per horizon, a filer's sample at a single horizon is
roughly its filing count. A per-firm HAR would then spend four parameters on a
couple of dozen noisy observations, where the pooled fit spends the same four on
the whole channel. Thin firm-level coverage shows up elsewhere in the evidence:
only 62.9 per cent of firms have any validation rows, which is why the
firm-identity reference falls back to a global validation mean for the rest. Those
covered firms do carry 91.5 per cent of long-form and 91.9 per cent of
event-driven test rows, so the sparsity is in the firm count and not in the
evaluation weight. % src: results/tables/firm_identity_control.md (firms with val rows; test-row coverage)
The per-firm arms that were executed carry fallbacks for the same reason, A3 and
A4 requiring 100 observations of history and A5 requiring 60 before they will
estimate at all. % src: configs/models/A3_garch.yaml, A4_egarch.yaml, A5_arima.yaml; src/sp500vol/models/price/garch.py, arima.py fallback paths

Second, the per-firm price fits the design does contain lose to pooled A2 under
two of the three loss conventions. A3 to A5 are estimated firm by firm and are the
three price challengers of the 180-comparison standalone census, contributing 27
comparisons over nine disclosure-by-horizon panels. Under the census convention of squared error, none
of the 27 favours the per-firm challenger and 23 are significantly worse than
pooled A2 after Holm correction; under volatility-unit QLIKE, none favours it and
15 are significantly worse. % src: results/tables/variance_unit_standalone180.csv (A3--A5 rows: better_se_holm 0, worse_se_holm 23, better_qlike_vol_holm 0, worse_qlike_vol_holm 15)
Under variance-unit QLIKE, seven of the 27 do favour the challenger, every one a
GARCH-family arm, six of them A3 and the seventh A4 at long-form $h=5$. A3's
own long-form margin points the same way at the two shorter horizons and the
other way at $h{=}20$, and is significant at none. % src: results/tables/variance_unit_standalone180.csv (better_qlike_var_holm; A3 long_form dm_qlike_var_clu -1.685/-0.792/+0.182 with holm 0.185/0.858/1.000)
Firm-by-firm estimation of a price model therefore buys nothing dependable on
this panel: it wins under one of three loss conventions and inside one model
family.

Third, the firm effect that pooling omits is supplied to the reference in
restricted form, and the consequence is measured. The firm-identity rung refits
the reference on the constant, the log HAR forecast and the log of each firm's own
mean validation realised volatility. That is one extra parameter whose coefficient
is shared across firms, a restricted firm effect rather than a free intercept per
firm, and it enters the reference rather than the text model. Against it the 38 of
69 cells that count as genuine under the single recalibrated HAR fall to 8 of 69,
and 38 of the 45 long-form cells turn negative. % src: results/tables/m1_ensemble_primary.md (38 of 69 genuine); results/tables/firm_identity_ensemble.md (8 of 69; 38 of 45 long-form negative)
A per-firm HAR would give the reference more than this rung does, at the cost of
the degrees of freedom just counted. What the rung establishes is that the cheap
version of that information removes most of the cells that looked incremental.

Fourth, the strongest rung of the ladder already contains per-firm price fits. The
maximal price reference is a validation-fitted log-space pool of five price
models, three of which (GARCH, EGARCH and ARIMA) are the firm-by-firm arms of this
block. Against it 9 of 69 cells survive Holm. % src: results/tables/maximal_reference_ensemble.md
The reference interval reported in Chapter~\ref{ch:results} is therefore bracketed
by references that do carry firm-specific price dynamics, whatever the HAR member
of the pool carries.

Two properties of the wider design bear on the same objection. The reference is
recalibrated on the validation years before any text term is added, at a mean
slope of 1.10 across the six channel-horizon cells. % src: results/tables/recal_slope_logspace.md
% (log-space b of Eq. 1; the level-space slope FAMILY 3 reports on the same rows
% averages 1.44, mean of 1.508/1.447/1.269/1.495/1.562/1.334)
The pooled fit's systematic under-response is therefore corrected rather than
left for text to repair.
And the direction of any residual weakness is known: a weaker price reference
credits text with more, so whatever pooling costs the reference inflates the most
generous rung's count of apparent increments rather than suppressing it. The
study's verdict rests on the stronger rungs, where the reference carries both the
firm term and the per-firm price fits.

\paragraph{Block~B.} B1 uses bag-of-words counts over 5{,}000 features with a
lower-case alphabetic token pattern of length $\geq 2$; B2 uses TF--IDF over
1--2-grams, 5{,}000 features, sublinear term-frequency scaling and the same token
pattern; B3 uses Loughran--McDonald Master Dictionary category proportions; B4
adds four engineered statistics (log token count, mean word length, mean sentence
length, punctuation density). % src: configs/models/B1..B4.yaml
All four fit ridge regression per horizon on the log target, with the penalty
chosen by five-fold cross-validation on mean fold squared error over
$\{0.1, 1, 10, 10^2, 10^3, 10^4\}$ for the two vectoriser arms and
$\{10^{-3}, 10^{-2}, 0.1, 1, 10, 10^2\}$ for the low-dimensional dictionary arms.
The folds carry a fixed shuffle state, so selection is deterministic. % src: src/sp500vol/models/classical_text/{bow_ridge,lm_linear,_fit_utils}.py

\paragraph{Block~C encoders.} S1 truncates to the first 512 tokens. S2--S4 split
the document into 512-token chunks with stride 256, keeping at most 8 chunks, and
run at physical batch 8 with 16 accumulation steps. S3 adds attention pooling of
width 128 over the chunk embeddings and S4 a one-layer, eight-head chunk-level
transformer encoder above them. % src: configs/models/C2_finbert_s{2,3,4}.yaml
C4 reads 4{,}096 tokens natively at physical batch 4 with 32 accumulation steps
and pre-tokenises in batches of 128, because tokenising 4{,}096-token windows of
$\sim$29{,}000-token documents on the fly would starve the card. % src: configs/models/C4_longformer.yaml;
% frozen 03_dataset.tex Table 1 (long-form median 29,278 wordpiece tokens)

\paragraph{C5 and C5x.} The three frozen embedders encode at maximum length
4{,}096 in encode batches of 8, L2-normalising the pooled embedding, under one
shared instruction: \emph{``Represent this corporate financial disclosure for
forecasting the issuing firm's future stock-return volatility''}. % src: configs/models/C5_*.yaml
C5x is the input-parity control: the same Qwen3-Embedding-8B encoder over the
byte-identical curated excerpts C6 reads, last-token pooling with L2
normalisation, and a ridge head on the log target with Duan smearing. It runs on
long-form only, as a single run. % src: release/run_configs/C5x_qwen3exc_full_long_form_seed2026.json; scripts/analysis/c5x_input_parity.py

\paragraph{C6 prompted decoding.} Qwen3-32B is served offline under vLLM in
bfloat16: temperature 0, thinking mode explicitly disabled (otherwise the output
budget is consumed by the reasoning trace), 8{,}192-token context with prompts
capped at 6{,}000 tokens, at most 120 new tokens, and guided JSON decoding against
the schema \verb|{vol_5d, vol_10d, vol_20d}| where the installed engine exposes
it, degrading to regular-expression parsing otherwise. % src: scripts/experiments/e1_llm_forecast/run_inference.py; TECHNICAL_SUPPLEMENT.md S5
A parse failure triggers exactly one retry pass carrying an explicit re-format
instruction at temperature 0.2. One call per filing returns all three horizons.
The \emph{curated excerpt} is defined mechanically: for 8-K filings the full text,
truncated only if it exceeds the $\sim$21{,}200-character excerpt budget. For
10-K filings it is Items 1A, 7 and 7A from the parsed section index and for 10-Q
filings Part~I Item~2 with Part~II Item~1A, used only when the concatenated
sections exceed 2{,}000 characters (below which the index holds
table-of-contents stubs). Otherwise the full document is head-truncated. % src: scripts/experiments/e1_llm_forecast/prompt.py
Forecasts are clipped to $[0.03, 3.00]$ annualised and a missing decode falls back
to trailing 22-day realised volatility. The shipped event-driven run records
39{,}322 filings and 117{,}407 rows at a 0.00\% parse-failure and 0.00\% clipping
rate. % src: release/run_configs/C6_llmtext_full_event_driven_seed2026.json stats block
The prompted arms train nothing, so only validation and test filings are ever
sent to the model.

\paragraph{Block~D.} D1 and D2 project price features and the FinBERT
\texttt{[CLS]} embedding to width 128 and then either concatenate them or combine
them through a learned gate, one weight per projected dimension rather than one
scalar; both read S1 512-token text at physical batch
128 without accumulation. % src: configs/models/D1_concat_mlp.yaml, D2_gated_fusion.yaml
D3 substitutes each frozen C5 embedding for the FinBERT branch under the
head-only recipe. D4 places the three HAR lags inside the C6 prompt and decodes
identically; its clipping rate is 0.07\% (88 of 117{,}407 rows). % src: release/run_configs/D4_llmfused_full_event_driven_seed2026.json


\begin{figure}[tbp]
\centering
\includegraphics[width=\textwidth,height=0.83\textheight,keepaspectratio]{figures/F2_model_spectrum_and_compute__a.pdf}
% caption numbers src: results/tables/seed_aggregate.csv (panel a, long_form slice: r2_mean, r2_std, n_seeds); results/tables/compare_full_long_form_test.md (the A2 HAR-RV R-squared triple, absent from seed_aggregate.csv); results/tables/cost_accuracy.csv (panel b, all 25 arms)
\caption[The model spectrum, and what the fine-tuning budget bought]{
The model
spectrum, and what the fine-tuning budget bought. 
Panel~(a): long-form
standalone $R^2$ for all fifteen text-only challengers at $h=5$, 10 and 20
(45 arm-by-horizon cells), as the mean over the seeds trained. Whiskers are at
$\pm$1 across-seed standard deviation where three seeds exist, open markers
where a single deterministic fit was made. Every one of the 45 lies left of
zero, from $-0.0149$ to $-0.4682$. % src: results/tables/seed_aggregate.csv (long_form slice, r2_mean over the 15 text arms x 3 horizons; extremes C5_qwen3 h=20 and B1_bow_ridge h=5)
Below the rule sit the five single price arms carrying a committed long-form
$R^2$ row, of which HAR-RV alone is positive at every horizon
($+0.0890/+0.1210/+0.2080$). Those five are \emph{not} the five members of the
maximal price pool; the panel says so. % src: results/tables/compare_full_long_form_test.md (A2 row); results/tables/seed_aggregate.csv (A1, A3, A4, A5 long_form rows)
The HAR-RV triple is the comparison table's reading, which predates the
smearing retransformation; the paragraph below sets it beside the committed
fit, whose long-form $R^{2}$ is higher at every horizon. % src: git 16c9d69^ against results/runs/A2_har_rv_full_long_form_seed2026/metrics.json (r2 0.0890/0.1210/0.2080 against 0.1694/0.1808/0.2640)
}
\label{fig:app-model-spectrum}
\end{figure}

\begin{figure}[tbp]
\centering
\includegraphics[width=\textwidth,height=0.83\textheight,keepaspectratio]{figures/F2_model_spectrum_and_compute__b.pdf}
% caption numbers src: results/tables/seed_aggregate.csv (panel a, long_form slice: r2_mean, r2_std, n_seeds); results/tables/compare_full_long_form_test.md (the A2 HAR-RV R-squared triple, absent from seed_aggregate.csv); results/tables/cost_accuracy.csv (panel b, all 25 arms)
\caption[The model spectrum, and what the fine-tuning budget bought (b)]{Continued from Figure~\ref{fig:app-model-spectrum}. The model
spectrum, and what the fine-tuning budget bought. 
Panel~(b): the same spectrum against recorded training compute, GPU-hours on a
broken logarithmic axis with the nine zero-cost arms in their own band,
against best long-form test QLIKE. Basis: means of the per-seed metrics over seeds
2026--2028 for every trained encoder and fusion arm, single deterministic fits
for the bag-of-words and price arms; the denominator of panel~(a) is 45 text
cells, the price block being drawn separately.
}
\label{fig:app-model-spectrum-b}
\end{figure}

Figures~\ref{fig:app-model-spectrum} and~\ref{fig:app-model-spectrum-b} make two points that do not survive
compression into a leaderboard. The first is that the whiskers of panel~(a)
are the reason the primary basis of this report is a seed ensemble rather than
any single run. At the frozen E5-Mistral arm at $h=20$ the across-seed
standard deviation is $0.2346$ against a mean of $-0.1369$, so several
whiskers cross zero even though no mean does. % src: results/tables/seed_aggregate.csv (long_form C5_e5mistral h=20: r2_mean, r2_std)
The second is that the panel must not be read as ``the price baseline wins''.
Several price arms sit below several text arms on $R^2$ while sitting above
them on QLIKE (ARIMA reads $-0.0816/+0.0206/+0.0666$, GARCH
$-0.2256/-0.1374/-0.2316$, EGARCH $-0.3057/-0.2156/-0.2306$ and historical
volatility $-1.3398/-0.7557/-0.3235$). % src: results/tables/seed_aggregate.csv (A1, A3, A4, A5 long_form r2_mean by horizon)
HAR-RV, the reference the whole
study is built against, is not the QLIKE-strongest price model, semivariance
HAR being better in all nine channel-horizon panels (variance-unit QLIKE-DM $-3.7$ to $-13.5$, observation-level and uncorrected, off the primary basis on all three counts). % src: results/tables/stronger_baselines.md (DM_q -3.735 to -13.504; "DM on the common (kept) test sample", no clustering and no Holm column)
That is what motivates closing the reference interval with a pool of price
models rather than nominating one, and it is the subject of
Section~\ref{app:res-frontier}.

What the figure does not show is incremental value. A poor standalone
forecaster can still remove loss beyond a recalibrated reference, which is the
object of the 69-cell combination grid and of neither panel here; no bar or
marker in Figures~\ref{fig:app-model-spectrum} and~\ref{fig:app-model-spectrum-b} entitles any cell to a survivor
count. The two panels also carry different quantities on different scales and
share no axis. The horizontal position inside the zero-GPU band is a
de-collision offset carrying no information, every arm in that band standing
at zero recorded GPU-hours with its CPU wall-clock seconds printed beneath.
The unit of panel~(b)'s QLIKE axis is not declared in its own source file,
which is why the axis is labelled as recorded. That last point is worth
closing, because the two committed tables feeding the two panels disagree
about the reference arm: the cost artefact records long-form test QLIKE
$0.5564/0.4618/0.2970$ while the standalone comparison table records
$0.7439/0.5794/0.3577$. % src: results/tables/cost_accuracy.csv (A2_har_rv row); results/tables/compare_full_long_form_test.md (A2_har_rv row)
Both vintages are printed on the artefact and no comparison is made across the
panels. The first of the two is the \emph{raw} HAR-RV forecast's variance-unit
QLIKE, which the combination grid records to six decimal places as $0.556448$
at long-form $h=5$; the recalibrated reference on the identical 7{,}951 rows
reads $0.510892$, so the vintage that matches is the uncombined one. The
second vintage is not that forecast under a different QLIKE convention: on the
same 7{,}951 rows it records MAE 0.1071, RMSE 0.1610 and $R^{2}$ 0.0890 where
the run's own committed metrics read 0.1058, 0.1537 and 0.1694, so it differs
on every metric and not only on the loss parameterisation. What separates the
two is one retransformation constant, and the repository settles it. All twelve
long-form figures of the comparison table reproduce, to every printed digit,
the metrics of an A2 run directory the run store no longer holds: nineteen
pre-rerun price and classical-text directories were removed on 2026-06-20, and
the removed A2 long-form directory is stamped 2026-06-08 where the committed
one is stamped 2026-06-22. % src: git 16c9d69 (2026-06-20, removal of nineteen pre-rerun baseline run directories); git 16c9d69^:results/runs/A2_har_rv_full_long_form_seed2026/metrics.json (test mae 0.107125/0.092738/0.077786, rmse 0.160975/0.156841/0.116420, r2 0.089036/0.120980/0.207980, qlike 0.743903/0.579353/0.357670) and env.json (2026-06-08T06:32:29Z) against the committed env.json (2026-06-22T18:28:32Z)
The two prediction files carry the same 94{,}237 rows, and joined row by row
their accession, ticker, effective trading day, split label,
realised-volatility label and all three HAR features differ by exactly zero, so
neither the sample, the split boundaries, the features nor the label
moved. % src: results/runs/A2_har_rv_full_long_form_seed2026/predictions.parquet against git 16c9d69^ of the same path (94,237 rows joined on accession-horizon-split; maximum absolute difference exactly 0.0 on label\_realised\_vol and feature\_rv\_1d/5d/22d; both carry train 2010-02-04--2019-12-23, val 2020-01-03--2021-12-23, test 2022-01-05--2025-12-23)
Their forecasts differ by a single multiplicative constant per horizon, steady
across all 94{,}237 rows to within $7\times10^{-12}$: $1.145059$, $1.100559$
and $1.076655$ at $h=5$, 10 and 20. % src: results/runs/A2_har_rv_full_long_form_seed2026/predictions.parquet against git 16c9d69^ of the same path (ratio of prediction\_realised\_vol, committed over removed, by horizon: minimum and maximum 1.1450585053373699/1.14505850534382, 1.1005586502880067/1.1005586502918845, 1.0766552112791599/1.07665521128133)
A constant ratio leaves the fitted slopes untouched, and those three numbers
are the smearing factors \cite{duan1983smearing}, recovered to eleven decimal
places as the mean ratio of label to forecast over the earlier fit's own
19{,}559, 19{,}543 and 19{,}500 training rows. % src: git 16c9d69^:results/runs/A2_har_rv_full_long_form_seed2026/predictions.parquet (split=train, mean of label\_realised\_vol over prediction\_realised\_vol by horizon: 1.1450585053379911, 1.1005586502884068, 1.0766552112794519)
The comparison table therefore records this same pooled fit on this same data
\emph{before} the smearing retransformation
specified above was written; the committed run is
stamped four minutes before the commit that adds it. % src: git a5feeff (2026-06-22 19:32:58 +0100, "Duan smearing retransformation for A2 HAR-RV baseline"), which also regenerated seven evidence tables including the DM and forecast-combination families

Three competing explanations are ruled out rather than left open. The loss
convention is not the difference: each vintage's predictions, passed through
the single metric function the repository has carried unchanged since its first
commit, reproduce that vintage's own committed MAE, RMSE, $R^{2}$ and QLIKE,
and that function squares its volatility inputs, so both readings are
variance-unit even though neither source table says so. % src: src/sp500vol/evaluation/metrics.py (all\_metrics squares before qlike; one commit on the path); both predictions.parquet vintages recomputed through it
The seed is not the difference, both configurations recording seed 2026 over a
deterministic least-squares fit. The pooling flag is not the difference either,
although the two configurations do disagree about it: the earlier declares
\texttt{fit\_per\_firm: true} and the current \texttt{false}, but no file under
\texttt{src/} or \texttt{scripts/} reads that field, and the earlier run's saved
model holds one coefficient vector per horizon and none per firm, so the flag
recorded a description that was wrong rather than an estimator that
changed. % src: config.json of both vintages; \texttt{fit\_per\_firm} occurs in configs/models/A2\_har\_rv.yaml and in no source file; git 16c9d69^:results/runs/A2\_har\_rv\_full\_long\_form\_seed2026/model.pkl (coefficients keyed 5/10/20 only)

Two consequences follow, and both are stated because the reference is what
every one of the 180 comparisons is measured against. The first is that the
correction runs against this report's own interest: retransformation raises the
forecast level, so the committed fit is the better of the two on every metric
--- long-form test $R^{2}$ rises from $0.0890$ to $0.1694$ --- and the price
reference the results must clear is therefore the harder one. The second is
that the superseded figures are confined. Eight committed files carry them,
every one a comparison table, and the only file in the code base that reads any
of them is the generator of Figures~\ref{fig:app-model-spectrum}
and~\ref{fig:app-model-spectrum-b}; the tables carrying the pairwise tests, the
combination grid and the identity ladder were all regenerated after the
correction and are dated after it. % src: results/ and release/ searched for A2\_har\_rv rows carrying the pre-correction QLIKE values (8 files: compare\_full\_\{long\_form,event\_driven,combined\}\_test.\{md,txt\}, compare\_runs\_full\_long\_form\_test.txt, release/aggregate\_results/compare\_full\_\{long\_form,event\_driven\}\_test.md); scripts/analysis/diss\_appendix\_figs/F2\_model\_spectrum\_and\_compute\_diss.py; last-commit dates 2026-06-16 for the comparison table against 2026-07-04 for dm\_pairwise\_clustered.csv and m1\_ensemble\_primary.csv
What the two panels disagree about is thus a dated correction, not an
unexplained discrepancy; each artefact is still printed as recorded, and no
number crosses between them.
The undeclared unit is why
Figures~\ref{fig:app-compute-accuracy} and~\ref{fig:app-compute-accuracy-b} label an axis
Figures~\ref{fig:app-model-spectrum} and~\ref{fig:app-model-spectrum-b} deliberately leave unlabelled.

\begin{figure}[tbp]
\centering
\includegraphics[width=\textwidth,height=0.83\textheight,keepaspectratio]{figures/FP4_compute_accuracy__a.pdf}
% caption numbers src: results/tables/cost_accuracy.csv (all 25 arms: gpu_hours_total, per-horizon long-form QLIKE, best_qlike, on_pareto_frontier); results/tables/cost_accuracy.md (block totals); Spearman computed in the generator over the 16 GPU-consuming arms
\caption[The efficiency frontier and the GPU-hour budget]{
The efficiency
frontier and the GPU-hour budget. 
Panel~(a) plots each of the 25 arms of the
efficiency artefact as its best long-form test QLIKE against the GPU-hours its
runs consumed, on a symmetric-logarithmic axis so the nine CPU arms can be
drawn at exactly zero. The Pareto frontier is a single point (GARCH, 0.2686
QLIKE at zero GPU-hours), and among the sixteen arms that did spend GPU time
the Spearman rank correlation between hours spent and loss is $+0.76$
($p = 0.0006$): the more expensive arms are the \emph{worse} ones. % src: results/tables/cost_accuracy.csv (on_pareto_frontier; Spearman over the 16 arms with gpu_hours_total > 0)
}
\label{fig:app-compute-accuracy}
\end{figure}

\begin{figure}[tbp]
\centering
\includegraphics[width=\textwidth,height=0.83\textheight,keepaspectratio]{figures/FP4_compute_accuracy__b.pdf}
% caption numbers src: results/tables/cost_accuracy.csv (all 25 arms: gpu_hours_total, per-horizon long-form QLIKE, best_qlike, on_pareto_frontier); results/tables/cost_accuracy.md (block totals); Spearman computed in the generator over the 16 GPU-consuming arms
\caption[The efficiency frontier and the GPU-hour budget (b)]{Continued from Figure~\ref{fig:app-compute-accuracy}. The efficiency
frontier and the GPU-hour budget. 
Panel~(b) shows where the 590.4 GPU-hours went, arm by arm on a logarithmic
axis, each row labelled with that arm's accuracy rank out of 25. The single
long-context arm consumes 254.7 hours, 43.1 per cent of the training budget,
and ranks sixteenth, while each of the three price-plus-embedding fusions
consumes 0.2 hours and ranks in the top eight. % src: results/tables/cost_accuracy.csv (C4_longformer 254.715 h at rank 16; D3 arms 0.20-0.21 h at ranks 6-8); results/tables/cost_accuracy.md (566.7 h for block C, 23.6 for block D, 590.4 combined)
Basis: cost is total GPU-hours over all runs of an arm (three seeds by
three disclosure channels for the encoder and fusion blocks). Accuracy
is the minimum over $h=5$, 10 and 20 of long-form test QLIKE in \emph{variance}
units, seed-averaged per horizon before the minimum is taken. Denominators:
25 arms in panel~(a), 16 for the rank correlation. The rank correlation's $p$
is a single uncorrected test of association across arms; the efficiency
artefact carries no pre-declared Holm family, and the correlation prices
architectures rather than an effect of compute.
}
\label{fig:app-compute-accuracy-b}
\end{figure}

Figures~\ref{fig:app-compute-accuracy} and~\ref{fig:app-compute-accuracy-b} and panel~(b) of
Figure~\ref{fig:app-model-spectrum-b} read the same 25 rows of the same
committed file, and are drawn twice on purpose. The earlier panel is an
orientation: it separates the zero-cost band, prints the CPU wall-clock
seconds those arms consumed, and states the unit ambiguity in its own source.
The later figure is the compute question at full resolution: it adds the
accuracy rank of every arm, the share of the budget each consumed, and the
rank correlation that turns ``the frontier is free'' from an impression into a
statistic. A reader who wants one compute figure should take
Figure~\ref{fig:app-compute-accuracy}; a reader who wants the spectrum in one
frame should take Figure~\ref{fig:app-model-spectrum}.

Four things the figure does not show bound the reading. It does not price
energy or carbon: no facility-level measurement was taken and none is
inferred, so 590.4 GPU-hours is an accounting of machine time and not an
environmental figure. It excludes prompted large-language-model inference
entirely, which has no row in the cost artefact, so the total is a
training-fleet figure and not a study-wide one. The two prompted arms,
one of which carries the study's only surviving residual, are therefore absent
from both panels. And the rank correlation is descriptive rather than causal:
the expensive arms are the fine-tuned encoders and the cheap ones the
frozen-embedding fusions, so the panel prices architectures and says nothing
about what a larger budget on a fixed architecture would do. Nor is the
accuracy ranking a ranking of usefulness: an arm can rank sixteenth on
standalone loss and still be one of the cells that adds at the first rung of
the reference ladder, which is the dissociation Section~\ref{app:res-itemgeography}
finds again inside a single document.

% No \clearpage before this section, and none before Section E.5 either. All
% seven section headings once carried one; five still do, and the five that
% remain were each tested and each earns its place. Every float in this chapter
% is set [p], so a \clearpage at a section boundary is what holds the section's
% figures inside it: remove the one before Section E.4 and Figure E.10, which
% belongs to Section E.3, lands after the E.4 heading, and remove the one
% before Section E.6 and Figure E.16 does the same to Section E.5. Removing the
% one before Section E.3 saves no page and opens a 78 per cent white one;
% removing the one before Section E.1 saves no page and takes the page that
% closes Section E.2 from 55 to 73 per cent white; removing the one before
% Section E.7 saves no page either. The two dropped here were the only two
% doing none of that work. Each left its section's closing paragraph alone on a
% page that was 85 and 80 per cent white, and dropping them takes this appendix
% from 40 pages to 38, clears those two pages and a third at 48 per cent, cuts
% a fourth from 73 to 55, moves no figure out of the section that discusses it
% and overflows no float. Measured by typesetting this file on its own at
% report geometry, one variant per \clearpage, not guessed. The arithmetic is
% what decides this, so re-measure before adding or removing a float here.

\FloatBarrier
\section{The Prompted Arm's Prompt, in Full}
\label{app:hyperparams:prompt}

Section~\ref{app:hyperparams:blocks} gives the decoding settings and the rule
that builds the curated excerpt. The strings themselves are printed here, so
that the instrument can be read rather than described. They are transcribed
from \texttt{scripts/experiments/e1\_llm\_forecast/prompt.py}, where the system
string, the task string and the message builder sit at lines 53, 59 and 137. % src: scripts/experiments/e1_llm_forecast/prompt.py
The transcription is verbatim, so the American spellings and the capitalisation
below are the source's and not this report's.

Every prompted call sends two chat messages. The system message is one string
and is identical in every variant.

\begin{verbatim}
You are a quantitative analyst who forecasts equity realised
  volatility from corporate SEC disclosures. You always answer
  with ONLY a single JSON object and no other text, no markdown,
  no explanation.
\end{verbatim}

Here and below, a line indented by two spaces is a typographic continuation;
every other line break is a newline in the string itself.

The user message is assembled from three pieces: a metadata header, a body and
a task instruction. The header is written first, by one code path, for every
variant.

\begin{verbatim}
Company SEC filing.
- Form type: 8-K (items: 2.02,9.01)
- Filing date: 2023-02-02
\end{verbatim}

The item list is the filing's \texttt{item\_subtype} field, and the parenthesis
is dropped when that field is empty, which it is for every 10-K and 10-Q: the
field is populated on all 117{,}407 event-driven rows and on none of the
35{,}635 long-form rows. % src: results/runs/C6_llmtext_full_{event_driven,long_form}_seed2026/predictions.parquet (item_subtype null share)
The header carries no ticker and no company name.

The body is what separates the arms. The full-text arm C6 inserts the curated
excerpt:

\begin{verbatim}
Filing excerpt:
<<<
[the curated excerpt]
>>>
\end{verbatim}

the date-only contamination control inserts one line in its place:

\begin{verbatim}
(No filing text is provided.)
\end{verbatim}

and the date-and-ticker control inserts that same line, preceded by a ticker
field appended directly to the header:

\begin{verbatim}
- Ticker: AAPL

(No filing text is provided.)
\end{verbatim}

The task instruction closes the user message, and the same string closes all
three of the arms above.

\begin{verbatim}
Task: based ONLY on the disclosure text above, forecast this
  company's ANNUALIZED realised stock-return volatility over the
  next 5, 10 and 20 trading days after the filing. Express each
  forecast as a decimal (e.g. 0.25 means 25% annualized; typical
  values lie between 0.10 and 1.00; extreme stress can exceed
  1.00).
Respond with ONLY this JSON object: {"vol_5d": <number>,
  "vol_10d": <number>, "vol_20d": <number>}
\end{verbatim}

Decoding is then constrained by the schema handed to the serving engine, where
that engine exposes guided decoding, which is why a malformed answer is rare:

\begin{verbatim}
{"type": "object",
 "properties": {"vol_5d":  {"type": "number"},
                "vol_10d": {"type": "number"},
                "vol_20d": {"type": "number"}},
 "required": ["vol_5d", "vol_10d", "vol_20d"],
 "additionalProperties": false}
\end{verbatim}

Two properties of this prompt read as holes on their own, so each is set out here with its answer.

\paragraph{The model is told the item code.} The header names the form type, the
filing date and, where the form carries them, the item codes, and it never
names the ticker or the company. The item codes are the part that needs saying out loud, because
Table~\ref{tab:app-item-qwen} partitions the 8-K residual by item code and so
cuts on a field the model was handed. The answer is in the composition printed
above. The builder writes the identical header and the identical task string
for the date-only control as for the full-text arm, and the two differ only in
the body, one carrying the document and the other a line saying there is none.
That control is significantly positive in 0 of 6
channel-horizon cells under Holm against the recalibrated HAR. % src: results/tables/llm_contamination.md (dateonly rows; Holm within the 18-cell block)
The metadata the header carries therefore buys nothing measurable on its own,
on either channel. The partition table is scored against a different reference and the
two sets of percentages are not interchangeable; what the control establishes
is the narrower point, that the item field is not itself the signal the
partition is cutting.

\paragraph{The task string states a prior.} Before the model has read anything
it is told that ``typical values lie between 0.10 and 1.00'' and that
``extreme stress can exceed 1.00''. That is an unconditional prior written into
the instrument rather than estimated from the training split, and a forecaster
that ignored the document and answered inside the band would already be roughly
scaled. The same control answers it. The task string is appended byte for byte
to the date-only and date-and-ticker controls as to the full-text arm, so both
hold the prior with no document at all. The date-only arm, which carries
nothing else either, is the one that scores 0 of 6 above; the date-and-ticker
arm adds the firm's name to the same prior and does not. The two elicitation paraphrases of
Chapter~\ref{ch:validation} restate the same band in different words, so they
vary the wording and not the prior. % src: scripts/experiments/e1_llm_forecast/prompt.py (_TASK_PARA1, _TASK_PARA2)
The bound is on the increment over the price reference, which is what this
report reads. It is not a claim that the prompted forecast is prior-free.

\FloatBarrier
\section{The Tuned Challenger Arm}
\label{app:hyperparams:asha}

The one arm that deliberately relaxes the fixed recipe is FinBERT under the S1
truncation strategy, specified in
\texttt{configs/hpo\_arm.yaml}, frozen at tag time.
% src: configs/hpo_arm.yaml tasks block (T1a and T1c both name model: C2_finbert_s1)
Sampling is Sobol
quasi-random with seed 2026 and adaptive search is banned. Asynchronous
successive halving promotes on pooled volatility-unit validation QLIKE through
rungs at 2, 5 and 15 epochs (2/5/10 for the Longformer task) with elimination
factor 3. The executed budget is \textbf{32} Sobol trials on the long-form
channel (task T1a) and \textbf{24} on the event-driven channel (T1c), the
event-driven search additionally running its first rung on a 30 per cent
subsample of training rows. % src: configs/hpo_arm.yaml sampler/asha/tasks (T1c carries rung0_subsample: 0.30; T1a does not)
The validation years are themselves split: fit rows to 2021-06-30 drive early
stopping and rung promotion, selection rows from 2021-07-01 rank configurations.
Test rows are \emph{physically deleted} before anything trains, and a SHA-256
manifest of the surviving row keys is stored per trial. % src: scripts/experiments/hpo/asha_hpo.py docstring and prepare_data
The eight-dimensional space covers learning rate (log-uniform
$5\times10^{-6}$--$1\times10^{-4}$), head learning-rate multiplier $\{1, 10\}$,
weight decay (log-uniform $10^{-4}$--0.3), head width $\{128, 256\}$, head dropout
$\{0, 0.1, 0.3\}$, freezing $\{$none, all but the top four encoder blocks, whole
encoder$\}$, effective
batch $\{32, 128\}$ and objective $\{$squared error, QLIKE$\}$; epochs is not a
dimension, being capped by the rung. % src: configs/hpo_arm.yaml space block (the three freeze_mode values none/top4/encoder); the semantics live in src/sp500vol/models/neural_text/bert_s1.py:442-451, where "top4" sets requires_grad=False on the whole encoder and then re-enables the last four transformer blocks, so it is the top four that stay trainable
The top two configurations by selection QLIKE were retrained under all three
seeds, the winner being the lower three-seed ensemble selection QLIKE.

The search selected trial 21 for long-form (learning rate
$9.11\times10^{-5}$, head multiplier 10, weight decay $1.57\times10^{-3}$, head
width 256, all but the top four blocks frozen, effective batch 32, QLIKE
objective) and trial 11 for event-driven (learning rate
$1.85\times10^{-5}$, head multiplier 10, weight decay $4.49\times10^{-4}$, head
width 256, dropout 0.3, no freezing, effective batch 128, QLIKE objective). % src: results/hpo/T1a/T1a_trial021_rung2/summary.json; results/hpo/T1c/T1c_trial011_rung2/summary.json
Both therefore depart from the fixed recipe in learning rate, head width and
training objective; Chapter~\ref{ch:results} reports what that bought.

\FloatBarrier
\section{Seeds, the Ensemble, and the Cross-Family Probes}
\label{app:hyperparams:seeds}

Every fine-tuned Block~C/D arm runs under seeds \textbf{2026, 2027 and 2028}. The
seed is applied before any data or model code executes, seeding Python's
\texttt{random}, NumPy and PyTorch \cite{paszke2019pytorch}; it is embedded in the
run identifier and in each per-horizon checkpoint fingerprint, so runs and
resumptions can never cross seeds. % src: configs/base.yaml; scripts/train.py
The \emph{seed-ensemble forecast} that serves as the primary evaluation object is
the per-observation arithmetic mean of the three seeds' predictions, formed by an
inner join on the filing-by-horizon key so that the ensemble is defined only on
rows every seed scored. % src: scripts/analysis/m1_ensemble_primary.py ensemble_text
The tuned arm of Section~\ref{app:hyperparams:asha} instead uses a per-observation
log-space mean, matching the rule that selected it; the two conventions are each
internally consistent and are never mixed.

Five cross-family replications re-read the identical prompt at temperature 0:
Llama-3.1-70B-Instruct in AWQ-INT4 quantisation (the \emph{only} quantised arm,
run tensor-parallel across two cards at 8{,}192-token context) together with
Yi-1.5-34B-Chat, Phi-4-14B, Mistral-Small-24B and Gemma-3-27B-Instruct, all in
bfloat16. % src: release/run_configs/C6_llmtext_{llama70,yi34,phi4,mistral24,gemma27}_*.json (llm field, quantisation); configs/prereg_residual_family_audit.md line 48 (vLLM offline batch, TP2, --max-model-len 8192, recorded as byte-identical to the C6/llama70 serving plan); scripts/experiments/e1_llm_forecast/run_inference.py:355 (--max-model-len default 8192); TECHNICAL_SUPPLEMENT.md Table S11
Three families (Llama-3.1-70B, Mistral-Small-24B and Gemma-3-27B) additionally carry three-run repeat ensembles. Because the runner
does not pass a seed to the serving engine and decodes at temperature 0, those
replicates differ only through kernel non-determinism: they are
reproducibility-jitter ensembles rather than stochastic-decoding ones, and
Chapter~\ref{ch:validation} reads them on that basis. % src: scripts/experiments/row15_llama70_ensemble/launch.sh header

The production fleet ran on NVIDIA A100-SXM4-40GB cards, four to a box with one
run per card. Physical batch sizes come from the per-arm configurations, which
set gradient accumulation so that every fine-tuned arm trains at the same
effective 128. A probe under an explicit headroom rule --- the largest passing
batch whose peak allocated memory stays at or below 36\,GB --- was run over two
arms only. The committed summary retains measured rows for those two
chunk-pooling arms (peak 29.1--29.2\,GB, first out-of-memory at batch 8 for
the hierarchical arm); the per-run configuration files, not the probe, are the
record of what executed. % src: results/tables/a100_40g_probe_summary.md
An earlier 24\,GB-class probe that sized a pilot campaign is retained as a tagged
pilot record and supports no production number. % src: results/tables/rtx3090_probe_24g_summary.md

\FloatBarrier
\section{Representative Run Commands}
\label{app:hyperparams:commands}

\begin{sloppypar}
All commands run from the repository root. Model identifiers are the
\texttt{model\_id} fields of \texttt{configs/models/}; \texttt{DISCLOSURE}
$\in \{$\texttt{long\_form}, \texttt{event\_driven}, \texttt{combined}$\}$.
\end{sloppypar}

\begin{verbatim}
# one fine-tuned arm, one seed (repeat over 2026 2027 2028)
python scripts/train.py --model C2_finbert_s1 --dataset full \
    --disclosure DISCLOSURE --seed 2027

# the production fleet: four cards, one run per card
python scripts/run_matrix.py --stage full --dataset full --gpus 4 \
    --seeds 2026 2027 2028

# per-run evaluation and cross-run grids
python scripts/evaluate.py --run-id C2_finbert_s1_full_long_form_seed2027 \
    --baseline A2_har_rv
python scripts/compare_runs.py --dataset full --disclosure DISCLOSURE \
    --split test --metrics mae rmse r2 qlike

# prompted arm: build the validation+test manifest, then decode offline
python scripts/experiments/e1_llm_forecast/run_inference.py build-manifest \
    --out results/e1_llm_forecast/manifest_valtest.parquet
python scripts/experiments/e1_llm_forecast/run_inference.py run \
    --manifest results/e1_llm_forecast/manifest_valtest.parquet \
    --model models/Qwen3-32B --variant both --max-model-len 8192 \
    --max-tokens 120 --out-dir /data/e1_raw_full --checkpoint-every 500
python scripts/experiments/e1_llm_forecast/postprocess.py build-runs \
    --raw-dir /data/e1_raw_full

# tuned arm: plan, search one rung across four shards, promote, select
python scripts/experiments/hpo/asha_hpo.py plan    --task T1a
python scripts/experiments/hpo/asha_hpo.py search  --task T1a --rung 0 \
    --shard 0 --num-shards 4
python scripts/experiments/hpo/asha_hpo.py promote --task T1a --rung 0
python scripts/experiments/hpo/asha_hpo.py select  --task T1a

# evidence layer: seed ensembles, the combination grid, the ladder
python scripts/analysis/m1_ensemble_primary.py
python scripts/analysis/forecast_combination.py
python scripts/analysis/firm_identity_control.py
python scripts/analysis/signal_injection_power.py
\end{verbatim}

Interrupted neural runs are resumed by re-issuing the identical
\texttt{train.py} command: completed horizons are recognised by their
fingerprinted checkpoints and skipped. Prompted decoding checkpoints every 500
generations --- one per filing and variant, not per filing --- so a preemption
costs at most the current chunk. % src: scripts/experiments/e1_llm_forecast/run_inference.py:249-253 and :276-281 (the work list holds one item per filing x variant and the flush counts those items)

\FloatBarrier
\section{Configuration Fingerprints}
\label{app:hyperparams:fingerprints}

Each run's configuration is canonicalised (serialised to JSON with sorted keys
and no separator whitespace) and hashed with SHA-256, so that a single changed
byte changes the digest. Across the fleet this gives \textbf{240} fingerprinted production runs
(smoke and sample probes excluded) carrying \textbf{233} distinct digests, with
\textbf{7} byte-identical pairs. % src: results/tables/config_fingerprints.md
Every one of those pairs is a single cross-family probe's combined and
event-driven run, byte-identical because that distinction is a downstream
row-subset selection rather than a configuration field: the fingerprint
identifies the configuration, not the run. % src: results/tables/config_fingerprints.md collisions table
A second digest hashes only the canonicalised training block, which is how the
one-recipe claim of Section~\ref{app:hyperparams:recipe} becomes checkable.

The released index publishes, per run, both the original digest and the digest of
the released preimage, together with the exact fields neutralised for anonymity
(38 of the 240 configurations named an absolute machine path, whose rewriting
necessarily changes the digest) and the 240 preimages themselves. Invoking
\texttt{config\_fingerprints.py} with the \texttt{-{}-verify-preimages} flag
re-hashes every released configuration against the published index. % src: results/tables/config_fingerprints.md method/release notes; release/config_hashes.csv header
A reader who wants to know exactly what produced a number in
Appendix~\ref{app:fullresults} can therefore go from the reported cell to a run
identifier, from the run identifier to a digest, and from the digest to the
configuration file that generated it.

\FloatBarrier
\section{Two Worked Filings, and the Prompt They Were Read With}
\label{app:vis-worked}

The object this section documents is a prompt and two documents rather than a
distribution, so it is set out in full rather than summarised into fields. Section~\ref{sec:res-residual} boxes the pair; what
follows is the same pair with the material the box had no room for, and the
exact prompt the arm was given.

\subsection*{The prompt, as issued}

The prompted arm receives two chat messages and returns one JSON object. The
variant shown is \texttt{c6\_text}, the text-only variant of the C6 arm; five
further variants share the template, among them the two paraphrases of
Table~\ref{tab:app-elicit} and the zero-content probes of
Figures~\ref{fig:app-elicitation-b} and~\ref{fig:app-elicitation-c}, which
delete the excerpt altogether. Reproduced verbatim, including
its American spellings and its emphasis capitals, the pair of messages for
Case~2 below reads:

\medskip
\noindent\fbox{\begin{minipage}{0.94\textwidth}
{\small\ttfamily\raggedright
\textrm{\textit{system message}}\\
You are a quantitative analyst who forecasts equity realised volatility from
corporate SEC disclosures. You always answer with ONLY a single JSON object
and no other text, no markdown, no explanation.\\[6pt]
\textrm{\textit{user message}}\\
Company SEC filing.\\
- Form type: 8-K (items: 8.01)\\
- Filing date: 2023-03-06\\[4pt]
Filing excerpt:\\
% T1 typewriter ligates << and >>, so the box printed «< and »> where the prompt
% holds <<< and >>>.  The braces break the ligature; the box claims to be verbatim
% and the paragraph below makes its character count checkable, so it must be.
<{}<{}<\\
\textrm{\textit{[the filing text, within the excerpt budget]}}\\
>{}>{}>\\[4pt]
Task: based ONLY on the disclosure text above, forecast this company's
ANNUALIZED realised stock-return volatility over the next 5, 10 and 20 trading
days after the filing. Express each forecast as a decimal (e.g. 0.25 means
25\% annualized; typical values lie between 0.10 and 1.00; extreme stress can
exceed 1.00).\\
Respond with ONLY this JSON object: \{"vol\_5d": <number>, "vol\_10d":
<number>, "vol\_20d": <number>\}
\par}
\end{minipage}}
\medskip

Of the three header lines only the latter two carry interpolated fields, so
those two and the excerpt are all that change from filing to filing. No
price, return or volatility quantity enters the text-only variant, which is
what makes it comparable with the trained text arms rather than with the
fused arm. The excerpt budget is 21{,}200 characters: a 6{,}000-token prompt
cap less a 700-token allowance for the scaffold, converted at the four
characters per token the prompt builder assumes off the GPU host, with the
host-side runner re-checking true token counts against the model's own context
limit. An 8-K shorter than the budget is passed whole and one longer is
truncated from the head, which is the policy that governs both cases here. % src: scripts/experiments/e1_llm_forecast/prompt.py (MAX\_PROMPT\_TOKENS, SCAFFOLD\_TOKENS, CHARS\_PER\_TOKEN, EXCERPT\_CHAR\_BUDGET, build\_excerpt)
That the template above is what the model actually saw is checkable rather
than asserted: rebuilding it over the two filings reproduces the released
prompt-length field exactly, at 21{,}723 characters for Case~1 and 4{,}789 for
Case~2. % src: release/raw_generations/raw/part-00197.parquet and part-00112.parquet (prompt\_chars, variant c6\_text)

\subsection*{Case 1: the document carries its own driver}

Fifth Third Bancorp, Form 8-K, items 1.01 and 9.01, accepted 2025-10-08 at
21:28:40 UTC and therefore assigned the effective trading day 2025-10-09;
accession 0001193125-25-234686, a test row. % src: release/accession_index.csv
The document reports entry into an agreement and plan of merger under which
Comerica is acquired for stock, each Comerica share converting into 1.8663
Fifth Third shares, conditional on both shareholder votes and regulatory
approval. % src: the filing itself, accession 0001193125-25-234686, indexed in release/accession\_index.csv. The released generation file carries the arm's output, its prompt length and an excerpt\_source flag taking the values full/head/sections; it does not ship the filing text, so the exchange ratio and the two conditions are verified against EDGAR rather than against that parquet
The document runs past the 21{,}200-character budget, so the arm
read the leading portion and not the tail. % src: release/raw_generations/raw/part-00197.parquet (excerpt\_source head, prompt\_chars 21723)
It answered \texttt{\{"vol\_5d": 0.35, "vol\_10d": 0.30, "vol\_20d": 0.25\}}
on the first attempt, with no retry and no parse failure. % src: release/raw_generations/raw/part-00197.parquet (raw\_output, parse\_ok true, retry\_used false)

Those three numbers sit at the top of the arm's own range. At $h{=}5$ the
value 0.35 is exceeded on 11 of the 25{,}109 event-driven test rows, and the
same row's price-only forecast from arm A2 sits at the 3.8th percentile of its
own distribution, the bottom decile at all three horizons. % src: results/runs/C6_llmtext_full_event_driven_seed2026/predictions.parquet; results/runs/A2_har_rv_full_event_driven_seed2026/predictions.parquet
The realised outcome fell in the top decile at $h{=}5$ and $h{=}10$ and the
eighth at $h{=}20$, and the text forecast carried the lower QLIKE at all three
horizons. % src: results/runs/C6_llmtext_full_event_driven_seed2026/predictions.parquet
One further reading of the same document is available and it is not
flattering. The fused D4 variant, given this excerpt together with three
trailing realised-volatility figures, answered 0.1500, 0.1650 and 0.1800, so
on the one filing of this pair where the text was right the addition of price
context removed what the text contributed. % src: release/raw_generations/raw/part-00197.parquet (variant d4\_fused, raw\_output)
That is one filing and settles nothing; the fused arm's standalone position is
the D4 row of Table~\ref{tab:res-leaderboard}.

\subsection*{Case 2: the driver is not in the document}

The same filer, Form 8-K, item 8.01, accepted 2023-03-06 at 21:30:14 UTC,
effective trading day 2023-03-07; accession 0000035527-23-000135, also a test
row. % src: release/accession_index.csv
The document announces that a previously disclosed accelerated share
repurchase of about \$200 million has finished: 5{,}589{,}996 shares
repurchased at an average price of \$35.78, leaving roughly 32.1 million
shares of authority outstanding. % src: the filing itself, accession 0000035527-23-000135, indexed in release/accession\_index.csv. As above, the released generation file records excerpt\_source = full (the whole document fitted the budget) but does not ship the text, so the three figures are verified against EDGAR
Read at face value it is a routine,
backward-looking capital-return notice, and if it points anywhere it points
away from turbulence. The whole text sits inside the
budget, so nothing was truncated. % src: release/raw_generations/raw/part-00112.parquet (excerpt\_source full, prompt\_chars 4789)
The arm answered \texttt{\{"vol\_5d": 0.18, "vol\_10d": 0.16, "vol\_20d":
0.14\}}, again first attempt. % src: release/raw_generations/raw/part-00112.parquet (raw\_output, parse\_ok true, retry\_used false)

The outcome fell in the top decile at all three horizons, above the 99th
percentile at $h{=}5$ and $h{=}10$. % src: results/runs/C6_llmtext_full_event_driven_seed2026/predictions.parquet
The price-only forecast missed it as well, sitting in the third decile, though
it carried the lower QLIKE of the two at all three horizons. % src: results/runs/A2_har_rv_full_event_driven_seed2026/predictions.parquet
The five trading days following 2023-03-07 lie inside the March 2023 stress in
United States regional banking, and the failure is therefore not a misreading
of the document: it is a demand for information the document does not contain.
No completed-buyback notice anticipates a sector-wide repricing three sessions
later, and the arm is given no price history with which to notice one
beginning.

That the miss is structural rather than particular can be shown on the same
committed rows. Filings whose effective trading days fall between 2023-03-06
and 2023-03-10 number 116, spread over 97 firms. A quarter of their $h{=}5$
rows land in the top decile of the realised outcome and 32.8 per cent in the
top two, against the tenth and the fifth that a decile cut allows. % src: results/runs/C6_llmtext_full_event_driven_seed2026/predictions.parquet
The arm's forecasts over the same window do not move with them: its median
forecast there is 0.18 against 0.22 across the panel, 1.7 per cent of its rows
carry a forecast of 0.32 or more against 2.4 per cent panel-wide, and it carries the
lower QLIKE against the price-only forecast on 24.1 per cent of those rows
against 43.1 per cent panel-wide. % src: results/runs/C6_llmtext_full_event_driven_seed2026/predictions.parquet; results/runs/A2_har_rv_full_event_driven_seed2026/predictions.parquet
This is the within-date decomposition of Section~\ref{sec:res-identity}
observed at the level of one week: the arm supplies cross-sectional ordering,
and a week in which the whole cross-section turns at once is where ordering is
worth least.

\subsection*{How coarse the arm's readings are}

The pair is legible partly because the arm's output space is small. Over the
25{,}109 event-driven test rows at $h{=}5$ it returns ten distinct values, and
two of them, 0.18 and 0.22, cover 92.6 per cent of the rows. % src: results/runs/C6_llmtext_full_event_driven_seed2026/predictions.parquet
Its Spearman correlation with the realised outcome is 0.17 against the
price-only forecast's 0.45, so on its own it is much the weaker ranking of the
two, which is the standalone result of Section~\ref{sec:res-standalone}
restated as a correlation. % src: results/runs/C6_llmtext_full_event_driven_seed2026/predictions.parquet; results/runs/A2_har_rv_full_event_driven_seed2026/predictions.parquet
What the arm does have is a usable high tail. On the 593 rows where it returns
0.32 or more, 36.1 per cent of outcomes land in the top two deciles against
20 per cent unconditionally. % src: results/runs/C6_llmtext_full_event_driven_seed2026/predictions.parquet
That contingency is reported here as description and is not an incremental
result: 43.0 per cent of those same 593 rows also sit in the price-only
forecast's top two deciles against 20 per cent unconditionally, so part of the
lift is information the price model already has, and conditioning on the
price-only decile shrinks the lift and reverses its sign in two of the ten
deciles. % src: results/runs/A2_har_rv_full_event_driven_seed2026/predictions.parquet; results/runs/C6_llmtext_full_event_driven_seed2026/predictions.parquet
No test is run on any of it.

\subsection*{How the pair was chosen, and what it is not}

The selection rule is stated so that it can be checked rather than trusted.
Case~1 was drawn from the 34 event-driven $h{=}5$ test rows on which the arm
returned 0.30 or more, the outcome fell in the top two deciles, the price-only
forecast fell in the lower six, and the text forecast carried the lower QLIKE. % src: results/runs/C6_llmtext_full_event_driven_seed2026/predictions.parquet; results/runs/A2_har_rv_full_event_driven_seed2026/predictions.parquet
Case~2 was drawn from the 84 rows on which the arm returned 0.18 or less, the
outcome fell in the top decile and the price-only forecast in the lower three. % src: results/runs/C6_llmtext_full_event_driven_seed2026/predictions.parquet; results/runs/A2_har_rv_full_event_driven_seed2026/predictions.parquet
Six filers appear in both pools, and one of the six was set aside because
its document concerns a named individual's alleged conduct, which
Appendix~\ref{app:ethics} treats as an incidental token in a model's input and
not as material to reproduce in a report. % src: results/runs/C6_llmtext_full_event_driven_seed2026/predictions.parquet (filers present in both pools)
The pair reported shares a filer, which is the property that removes firm
identity from the contrast between the two cases and is the reason this
particular pair was preferred over the others available.

Four things the pair does not establish. It was chosen for legibility and is
not a random draw from either pool, so its two outcomes carry no sampling
interpretation. It enters no multiplicity family, and no count, test or
interval in Chapters~\ref{ch:results} or~\ref{ch:validation} moves by a single
unit whether it is included or not. Every outcome above is a decile or a
percentile within the event-driven test distribution at the horizon named,
because the release withholds per-filing realised volatility and a level would
disclose exactly the quantity it withholds; the deciles are computed on the
committed prediction files rather than published as a table. % src: release/DATA_CARD.md; results/runs/C6_llmtext_full_event_driven_seed2026/predictions.parquet
And the elicited forecast shown is not the quantity the chapters test. Those
tests concern the combination of the text forecast with a recalibrated price
reference, and the combination enters the text forecast with an
out-of-sample log elasticity of $+0.264$ in this cell, so a swing from 0.18 to
0.35 reaches the combined forecast heavily damped. % src: results/tables/forecast_combination_grid.csv (event\_driven, C6\_llmtext, h=5, g\_log)
That damping is one reason an increment this visible on a single document
aggregates, in this cell, to $+1.21$ per cent over the recalibrated price
reference --- about one per cent, not a fraction of one, and smaller again once
the firm-identity term enters the reference. % src: results/tables/control_intersection_ensemble.csv (event_driven C6_llmtext h=5, vol_rel_impr_pct 1.214)

\clearpage
